She walked into the Great Hall on the second morning and found Samwell Tarly and said, “Take it off.”

He was elbow-deep in somebody’s thigh at the time and did not look up for a moment, and when he did his face went through about four expressions in rapid order before settling on the last one, which was relief.

“Oh,” he said. “Oh, thank the gods. I’ve been trying to work out for a day and a half how to make you let me look at it.”

“You’d have got nowhere.”

“I know. I asked Jon and he said not to bother trying.” Sam handed off what he was holding to a girl of about fifteen with blood to her wrists, gave her three instructions Arya did not follow, and wiped his hands on a rag that was not clean. “Sit down. Not there — there, by the light. Can you get the strapping off yourself or do you want me to cut it?”

“Cut it.”

He cut it, and the arm came away from her body and hung, and Arya looked at it properly for the first time since the middle of the first night.

It was not her arm. That was the plainest way to put it. It was grey-white to the elbow and slightly swollen and the skin had gone shiny and there was a black line running up the inside of the forearm from the wound, and there was a smell coming off it that she had smelled perhaps two hundred times in the last three days and had never once expected to be attached to herself.

“Right,” Sam said. He did not flinch and he did not perform not-flinching either, which she appreciated more than she expected to. “How much of it can you feel?”

“None of it below the elbow. There’s something at the elbow. Not pain. Pressure.”

“How long?”

“The elbow started yesterday. Around the time the fighting stopped.”

“Then it’s moving up.” He sat back on his heels. “Arya, I’m going to tell you exactly what I think and I’d rather do it badly and quickly than well and slowly. That arm has been dead since the first night and it’s begun poisoning the rest of you, and the black line means it’s already past where I want it to be. If I take it below the elbow I don’t think I’ll get all of it and we’ll be doing this again in four days and you’ll be considerably worse when we do. If I take it above the elbow I probably will.”

“Then above.”

“You want to think about it.”

“No.” She had thought about it. She had thought about it for the whole of the previous day, in a corridor, alone, with her back against a wall, and had arrived at the arithmetic somewhere around midnight and had spent the hours since being angry at herself for the twelve hours she had wasted first. “You’ve got other people. Do it now.”

Sam looked at her for a moment longer.

“All right,” he said.

\sceneBreak

There was no milk of the poppy.

“There’s about four doses left in the city,” Sam said, while the girl of about fifteen was laying things out on a board. “I’ve got them for men who might still die of the shock. You won’t. I’m sorry. That’s the whole of the reason and it’s not a good one but it’s the true one.”

“Fine.”

“It isn’t fine, and I’d rather you didn’t say it was, because in about four minutes you’re going to find out how not fine it is and I’ve had people hate me afterward for letting them say fine beforehand.”

“Sam.” She looked at him. “It’s fine. I’ve had worse from people who were enjoying it.”

He did not have an answer to that, and she had not really meant him to.

They gave her a leather strap for her teeth and four men to hold her, and she told them to use three because the fourth was needed elsewhere, and they used four anyway.

\sceneBreak

She was clear-eyed through the whole of it, which she had not expected and did not enjoy.

That was the thing nobody warned you about. She had been hurt a great many times and had learned that severe pain generally arrived with a kind of merciful narrowing — the world went small and loud and afterward there were gaps. This did not do that. She was awake and lucid and present for every single part of it, and what she found herself doing, with a detachment that horrified her later, was \textit{noticing}.

She noticed that Sam’s hands did not shake. Not once, not at any point, and she had watched that man’s hands shake while pouring wine.

She noticed that he talked the whole time — not to her; to the girl of fifteen, a low continuous stream of instruction in the completely level voice of a man reading aloud, \textit{now the retractor, no the other one, good, hold that, keep holding it, don’t look at me, look at what your hands are doing} — and that the girl held everything he told her to hold.

She noticed the exact moment the saw started because it was a different sound than everything before it, and she noticed that Sam did not slow down.

She noticed, at about the point where she thought she was going to pass out and then failed to, that somebody two pallets over was singing. Some Northern thing. Badly.

And she noticed, afterward, in the long grey time when they had finished and she was lying on her back with her teeth aching from the strap and the whole left side of her body reporting an enormous absence, that Samwell Tarly had gone about eight feet away, put both hands on a table, and stood there breathing for perhaps twenty seconds before going to the next person.

\textit{He hated that,} she thought, distantly. \textit{He hated every second of that and he did it anyway and he didn’t get it wrong.}

She had spent two years around people who killed with tremendous skill and total composure. She found, lying on a pallet in a converted throne hall with one arm, that she had never in her life been so impressed by anybody.

\sceneBreak

Jon came in the afternoon, or possibly the evening. She had lost the shape of the day.

He sat down on the floor beside the pallet rather than standing over her, which she noted, and he held his right hand in his left in his lap, wrapped to the wrist in linen, and Arya looked at it and did the sum.

“He saved it.”

“Aye.” Jon turned it slightly, without unwrapping it. “Most of it. Two of the fingers won’t ever do much. The rest — he says it’ll be stiff and it’ll shake when I’m tired, and it’ll ache in the cold for the rest of my life, and I’ll not know from one morning to the next whether it’ll hold a sword properly.” A pause. “He said that last part to me straight out. Nobody else would have.”

“That’s worse than losing it.”

“Aye,” Jon said. “I’ve thought that once or twice already.”

They sat with that for a while. There was a great deal of noise in the hall.

“He didn’t think he’d save it,” Jon said eventually. “I want you to know that, because he’ll not tell you himself. Sansa told me they’d all agreed it was coming off — she’d said as much out loud the first night. And he went in and worked on it for four hours and got the dagger out of the fist and cleaned what was underneath and then he came and sat down next to me and cried for about a minute and then told me he thought he’d got it.”

“Four hours.”

“Four hours. On one hand. With two hundred people in this hall who needed him.” Jon’s mouth did something complicated. “I asked him afterward why he’d spent four hours on one man’s hand with all that going on and he said — “ He stopped, and started again. “He said he’d spent his whole life being able to tell people exactly what was happening to them and never once being able to do anything about it. He said he’d wanted one.”

Arya looked at the ceiling of the Great Hall for a while.

“Well,” she said. “He got one.”

\sceneBreak

“Where’s the dagger?” she asked, later.

Jon did not answer immediately, which told her he had been waiting for the question and had not decided on the answer.

“Sam’s got it. He took it off me and he’s not given it back and I’ve not asked him to.”

“It’s mine.”

“I know it’s yours.”

“Bran gave it to me.” She heard her own voice do something she did not much like. “He gave it to me at Winterfell and he said — he didn’t say anything, actually. He just gave it to me. But it was a \textit{giving}, Jon, it was the only thing he ever did after the cave that was like a person doing something.”

“I know.”

“So it’s mine, and I lost it, and I want it back.”

Jon was quiet.

“Do you?” he said.

And Arya lay on a pallet with one arm and discovered that she did not have an answer, because the honest thing — the thing she had been circling for a day and a half and had not permitted herself to look at directly — was that she had let go of it on purpose. She had made a decision in a second and a half with a dead thing on her chest and the decision had been \textit{I need the hand}, and she had opened her fingers, and the world had been saved by that specific selfish choice, and she was never going to be able to tell anybody about it in a way that came out sounding true.

“I let go of it,” she said.

“Aye. Sam worked that out. It was in the fire before I got to it and there was nobody else out there who’d have been carrying it.”

“That’s not — “ She stopped. “I didn’t drop it and I didn’t have it knocked out of my hand. I opened my fingers, because there was a thing on top of me and I needed the hand, and I chose the hand.” She turned her head and looked at him. “You went into a fire and I let go of a knife to save my own life, and it’s the same knife, and everybody in this city is going to spend the next hundred years telling one of those stories.”

Jon looked at his wrapped hand for a long moment.

“I didn’t decide anything,” he said. “I want that said once, out loud, to somebody, and then I’ll not say it again. There wasn’t any deciding. I saw a blade in a fire and I’d watched steel do nothing for two days and my hand was in it before I’d finished the thought. If I’d had one more second I’d like to think I’d have done it anyway, but I didn’t have one, so I’ll never know, and neither will anybody else.” He looked up. “You had a second and a half and you did arithmetic. That’s harder, Arya. That’s considerably harder and it doesn’t make a song and I’d not swap with you.”

She did not say anything for a while.

“Keep it,” she said, finally. “The dagger. Or let Sam keep it, or throw it in the Blackwater with that horn, I don’t much care. I don’t want it back.”

“Arya — “

“I can’t hold it anyway.” She lifted the stump about two inches off the pallet, which was as far as it went, and it hurt in a way that had nothing to do with the arm. “That’s not self-pity, that’s a fact about a knife. I threw with the left. Every single time, my whole life — Needle in the right and the small one in the left, and I’ve been left-handed with a knife since Syrio Forel told me to stop being predictable.” She let it drop. “So somebody else can have it. I’m not going to be the woman who carries a Valyrian dagger about in one hand for the rest of her life so people can point.”

\sceneBreak

She thought about the other thing that night, when the hall had gone as quiet as it ever got.

Braavos was closed.

That was the piece she had not said to Jon and did not intend to say to anybody. The work she had been trained for — the actual work, the discipline, the thing she had given up two years and a great deal else to learn — required two hands. You could not take a face with one hand. You could not do the wearing of it, the fixing, any of it. She had known that within about four minutes of Sam telling her the arm was moving up and had made herself finish the arithmetic anyway.

She had spent six years knowing that if everything here went wrong, if the North fell, if her family died, if she ran out of names — there was a house across the sea with a door that would open for her.

The door was shut now. Not by anybody. By a wound.

And what she noticed, lying in the dark with the whole left side of herself gone strange, was that she could not tell whether the thing she felt about it was grief.

She was still faster than most people. She would still be faster than most people in a year, with one hand, because being fast had never actually been about hands. But the \textit{shape} of what she was for had been built around a very specific set of capabilities, and about a third of them had gone into a bucket that morning, and Arya Stark lay awake in a hall full of dying men and understood that she was going to have to work out what was left of her, from the beginning, at nineteen years old, with no list and nothing to kill.

She did not find that frightening, exactly.

She found that she had absolutely no idea how a person went about it.